\section{Short trajectories and local certificates}\label{lp:section}\label{sec:local-paths}

The energy estimates control occupations within cells.  We now prepare
the comparison of deletion measures needed to control variation between
cells.  This section constructs paths that can be resampled while
their endpoints remain fixed.  The required conclusions are quantitative:
the excluded mass must be smaller than $D^{-1}$, and every local insertion
must have a probability bounded away from zero, uniformly in the volume.
We prove these statements before using the paths to compare particle-deleted
configuration measures.  The stationary path representation and removal
argument come from~\cite[Section~4]{H}.
The short time scale and the finer spatial certificates below require
additional estimates.

Throughout this section $\Phi\geq0$ is a normalized bosonic ground function,
$\mu(\dd X)=\Phi(X)^2\dd X$, and $E_N$ is its eigenvalue.  We use the
scaled torus of side $K$, density $D=N/K^3$, and exclusion distance
$r=\alpha/D$, with $\alpha$ in the fixed compact window already chosen.
Fix a sufficiently small number $x>0$, to be chosen after the exponents
in the energy localization, and set
\begin{equation}\label{lp:scales}
 t=\frac1{x\log D},\qquad f_0=D^{-x},\qquad s=D^{-4/5},
 \qquad b=\frac{C_1}{\sqrt{x}},\qquad L_0=C_2\sqrt{\log D}.
\end{equation}
The numerical constants are chosen in the order $C_1$, then $C_2$.
Constants may depend on the fixed $x$ and on fixed neighborhood sizes,
but never on $K$, $D$, or $\Phi$.  A factor $D^{o_x(1)}$ means a bound
$C(x)D^{\zeta(x)}$ with $\zeta(x)\to0$ as $x\downarrow0$.
Such factors will occur only a bounded number of times in a local
comparison; this notation supplies no allowance per step of a long route.

\subsection{The stationary law and deletion of paths}

Let $\nu_{N,T}$ be the measure on $N$ independent torus Brownian motions
of generator $\Delta$ on $[-T,T]$, with Lebesgue initial measure.  Define
\begin{equation}\label{lp:path-law}
 \dd\mathcal R_T(\omega)
 =e^{2TE_N}\Phi(\omega(-T))\Phi(\omega(T))
  \1_{\{d_K(\omega_i(q),\omega_j(q))>r\ \forall q,\ i<j\}}
  \dd\nu_{N,T}(\omega).
\end{equation}
Here $d_K$ denotes torus distance.  Every time marginal of
$\mathcal R_T$ is $\mu$; restriction to a shorter interval gives the same
formula with its own endpoints.  Conditional on the endpoint
configurations, the law is the product torus-bridge law conditioned on
strict mutual avoidance.  Specifying each winding gives the corresponding
Euclidean bridges.  These facts follow by splitting the killed heat
semigroup and using $e^{-qH_N}\Phi=e^{-qE_N}\Phi$; they are the
stationary-path lemma of~\cite[Lemma~4.1]{H}.  In particular the law on $[-t,t]$
is the restriction of the law on $[-1,1]$ once $D$ is sufficiently large.

We will repeatedly delete the interactions involving a few specified
labels.  The useful version allows different weights on these labels.
If $J$ is a set of $m$ distinct labels and $H_i$ are nonnegative
single-path functionals invariant under time reversal, then
\begin{equation}\label{lp:deletion}
 \E_{\mathcal R_T}\prod_{i\in J}H_i(\omega_i)
 \leq e^{CT\alpha m}
       \int\prod_{i\in J}\E_{
          \mathrm{free}}^{Y_i}H_i\,\mu(\dd Y),\qquad 0<T\leq1.
\end{equation}
Here $Y$ is the initial configuration of those free paths, whose
duration is $2T$.  For completeness,
after dropping the selected interactions, the remaining killed
semigroup is bounded by $e^{-2TE_{N-m}}$ between the corresponding
slices of $\Phi$.  If $a(y)$ is the $L^2$ norm of such a slice, its
endpoint factor is $a(y)a(z)$.  The selected weighted endpoint kernel
is symmetric, so $a(y)a(z)\leq(a(y)^2+a(z)^2)/2$ replaces this factor
by $a(y)^2$ after integration.  The removal bound in
Proposition~\ref{set:removal}, $E_N-E_{N-m}\leq C\alpha m$,
then proves~\eqref{lp:deletion}.
Truncation proves the assertion for unbounded weights.  This is the
proof of the deletion lemma in~\cite[Lemma~4.2]{H}; using distinct reversal-invariant
weights does not change its symmetric-kernel argument.  A weight that
is not reversal invariant can always be bounded by its sum with its
time reversal.

We first use the longer slab to estimate close endpoint pairs.  For any
fixed bounded region $U$ and any $q\in[-t,t]$,
\begin{equation}\label{lp:close-pairs}
 \E\sum_{i\ne k}
  \1_{\{\omega_i(q)\in U,\ d_K(\omega_i(q),\omega_k(q))\leq s\}}
 \leq C_U D^2s^3.
\end{equation}
To prove this, cover $U$ by $O_U(s^{-3})$ cubes of side $s$ and enlarge
each cube by a fixed factor.  A pair counted on the left lies in one
of the enlarged cubes.  In the free law on $[-1,1]$, observation at
$q\in[-t,t]$ is uniformly separated from either slab endpoint.  The
heat-kernel estimate for a cube $Q$ therefore gives
\[
 \mathbb P_{
 \mathrm{free}}^y\{\omega(q)\in Q\}
 \leq Cs^3e^{-2d_K(y,v_Q)},
\]
where $v_Q$ is a nearby unit-cell center; changing its position by a
bounded amount changes only $C$.  Apply~\eqref{lp:deletion} to each
pair, and sum the resulting products.  The count moment-generating
bound in Proposition~\ref{set:count-mgf} implies
$\E(\sum_i e^{-d_K(Y_i,v_Q)})^2\leq CD^2$.
Thus one cube costs $CD^2s^6$, and the covering proves
\eqref{lp:close-pairs}.  The estimate holds uniformly under translations
of $U$ and does not assume a translation-invariant ground function.

\subsection{Exposure and the meaning of a certificate}

Write $Y_i=\omega_i(-t)$, $X_i=\omega_i(0)$, and
$Z_i=\omega_i(t)$.  Divide each half slab into equal intervals of
length $\Delta\in[r^2/2,r^2]$.  A path satisfies the increment cutoff if
its Euclidean lift obeys
\begin{equation}\label{lp:cutoffs}
 |\widetilde\omega(q)-\widetilde\omega(q')|
 \leq L_0\sqrt{|q-q'|+r^2}
 \qquad(q,q'\in[-t,t]).
\end{equation}
The lift increment is independent of the initial representative.
When $Y_i,Z_i$ have near representatives at distance at most $b$, let
$m_i$ be their Euclidean mean.  For sufficiently large $K$ the near
displacement is unique.  Call label $i$ eligible if both of its
endpoints are more than $s$ from every other corresponding endpoint,
its near endpoint displacement is at most $b$, and its path follows
that winding, satisfies~\eqref{lp:cutoffs}, and stays within $20b$ of
$m_i$.

Independently assign each label a Bernoulli flag of mean $f_0$.  Let
$\mathcal I$ be the set of flagged eligible labels.  Expose the endpoints,
flags, membership in $\mathcal I$, and the full paths of all labels outside
$\mathcal I$.  The latter are called frozen paths; the exposed variables
collectively form the datum $\mathcal D$.  An interior label at a site is a member
of $\mathcal I$ whose initial endpoint is in that unit cell and whose endpoint
displacement is at most one.  Tags used at the ends of comparisons
may have displacement up to $b$; their site is the unit cell containing
$m_i$.

Let $B_i$ denote the plain Euclidean bridge probability from $Y_i$ to
$Z_i$ in time $2t$ along the selected winding.  For a fixed datum put
$\mathcal A_i$ for the event consisting of label $i$'s individual
cutoffs and avoidance of all frozen paths.  The exact conditional law
of the unexposed paths is
\begin{equation}\label{lp:conditional-law}
 \lambda(\dd\omega_{\mathcal I})
 =\frac1{Z_{\mathcal I}}
    \prod_{i\in\mathcal I}\1_{\mathcal A_i}(\omega_i)
    \prod_{\substack{i<k\\i,k\in\mathcal I}}
      \1_{\{d_K(\omega_i(q),\omega_k(q))>r\ \forall q\}}
    \prod_{i\in\mathcal I} B_i(\dd\omega_i).
\end{equation}
This identity holds for almost every sampled datum, and then $Z_{\mathcal I}>0$.
Indeed, conditional on endpoints and flags, eligibility of each
unexposed label is an individual test on its own trajectory.  The
tests for the other labels become fixed when their paths are exposed.
Confinement chooses a single winding; its torus heat-kernel factor is
constant on the conditional fiber and cancels.  Disintegration of
\eqref{lp:path-law} gives~\eqref{lp:conditional-law}, exactly as in
\cite[Proposition~6.2]{H}.  No condition involving
the collection of latent trajectories is added to this law.

Choose a neighborhood radius $R_*=1000b$, enlarged by fixed factors
when necessary to contain all trials under discussion.  For a mesh
interval $J$ write
$A_J(\omega)=r^{-1}\operatorname{diam}\widetilde\omega(J)$.
We shall control collisions with unexposed paths by their endpoint
counts, and collisions with frozen paths by their counts and oscillations
at mesh times.  The following certificate records precisely this
information.  Every test uses exposed data only, so it leaves
\eqref{lp:conditional-law} unchanged on each conditional fiber.

\begin{definition}[Good site]\label{lp:certificate}
A unit-cell site is good if all of the following conditions hold.
\begin{enumerate}
\item It has at least $c f_0D$ interior labels.  In its $R_*$
neighborhood, each unit cell has at most $CD$ endpoints at either
time and at most $Cf_0D$ flagged endpoints at either time.
\item At either endpoint time, every ball $B(z,h)$ in the neighborhood,
with $r\leq h\leq1$, satisfies
\begin{equation}\label{lp:flagged-count}
 \#\{i:\ i\text{ flagged},\ Y_i\in B(z,h)\}
 \leq C(f_0Dh^3+D^{1/100}),
\end{equation}
and the same assertion holds for $Z_i$.
\item Every frozen path visiting the relevant neighborhood satisfies
\eqref{lp:cutoffs} and has lifted diameter at most $Cb$.  At each mesh
time $q$, and for either adjacent mesh interval $J$, the frozen paths
satisfy
\begin{equation}\label{lp:frozen-count}
 \sum_{k\notin\mathcal I:\ \omega_k(q)\in B(z,h)}
        (1+A_J(\omega_k))^3
 \leq C(Dh^3+D^{1/100})\qquad(r\leq h\leq1).
\end{equation}
The unweighted version is also required at endpoints and at time zero.
\end{enumerate}
All ball tests may be made on finite spatial nets with dyadic radii,
using enlarged balls and fixed margins.  The stated inequalities for
arbitrary balls then follow by covering.  The constants $c,C$ are
fixed before taking $D$ large.  Write $\mathcal B$ for the set of bad sites.
\end{definition}

For later probability estimates we also use the event obtained by
requiring the third item for the entire path collection, with the
same neighborhoods and fixed margins.  In addition, each original path
visiting the neighborhood is required to obey both ordinary-motion bounds
\begin{equation}\label{lp:ordinary-motion}
 \sup_{q\in[-t,t]}|\widetilde\omega(q)-\widetilde\omega(-t)|\le b,
 \qquad
 \sup_{q\in[-t,t]}|\widetilde\omega(q)-\widetilde\omega(t)|\le b.
\end{equation}
We call their intersection the full-path event.  It is an event in the
original sample, used to bound discarded mass; it is not a restriction
in~\eqref{lp:conditional-law}.  In particular, for an original midpoint
$y=\omega_k(0)$ on this event, $|y-m_k|\le b$.  If label $i$ has a
neighboring mean, $|m_i-m_k|\le C$, then for sufficiently small fixed
$x$ one has $|y-m_i|\le2b<8b$.  This is the midpoint range needed in
the pinned insertion estimate below.

\begin{proposition}[Probability of certificates]\label{lp:good-probability}
The constants in Definition~\ref{lp:certificate} can be chosen so that
\begin{equation}\label{lp:good-probability-bounds}
 \mathbb P\{v\text{ is bad}\}\leq C(x)D^{-1-1/8},
 \qquad
 \mathbb P\{A\text{ consists of bad sites}\}\leq p^{|A|}
\end{equation}
for every site $v$ and every set $A$ of distinct sites.  Here $p>0$
can be any prescribed fixed constant, followed by a sufficiently large
lower threshold on $D$.  Failure of the full-path event in a fixed
neighborhood has probability at most $C(x)D^{-6}$.  The estimates are
uniform in $\Phi$ and in all sufficiently large torus sizes $K$.
\end{proposition}

\begin{proof}
We give the separate estimates for supply, exceptional motion, and
the fine ball counts.  Their combination will also explain why the
one-site power and the joint parameter in
\eqref{lp:good-probability-bounds} have different forms.

First consider the cutoffs under free motion.  Gaussian maximal
estimates give probability at most $CD^{-40}$ for maximal displacement
from the initial point larger than $b$ or failure of~\eqref{lp:cutoffs},
after the choices of
$C_1,C_2$.  For the increment estimate, first use a union bound over
the $O(r^{-4})$ pairs of mesh points and then control oscillations on
individual intervals.  For a trial bridge with endpoint displacement
at most $b$, or two legs pinned at a point within $8b$ of $m_i$, the
increment cutoff and confinement to the radius-$20b$ ball fail with
probability at most $CD^{-40}$.  The maximal-displacement test at
radius $b$ is used only for exceptional visitors in the original
sample.  To check the trial estimates uniformly as $t$ decreases,
represent the centered
part of a leg of length $h$ by $W(q)-(q/h)W(h)$.  Its prescribed line
has increment divided by $\sqrt{|q-q'|+r^2}$ at most
$C b/\sqrt t=C C_1\sqrt{\log D}$, so it uses only a fixed fraction
of $L_0$.  The Gaussian remainder and the distance from that line to
the confinement boundary give the asserted tails.  Reversal of any
of these exceptional events preserves the bound.

Here is the transfer of these free tails to the interacting sample.
For a set $A$ of sites let $J_A(\omega)$ count the neighborhoods in
that set visited by a path.  If a reversal-invariant path event
$\mathcal E$ has free probability at most $\varepsilon$, the free
visit estimate in~\cite[Lemma~4.5]{H}, with a
unit-cell covering for our fixed neighborhood radius, gives
\[
 \E_{\mathrm{free}}^y
       [e^{\vartheta J_A\1_{\mathcal E}}-1]
 \leq C_\vartheta\varepsilon^{1/3}e^{-d_K(y,A)}.
\]
Its constants are uniform for $t\leq1$.  Expanding the exponential
as a product over labels, applying~\eqref{lp:deletion}, and then
using the count moment-generating estimate yields
\begin{equation}\label{lp:visitor-mgf}
 \log\E\exp\!\left(\vartheta\sum_i
             J_A(\omega_i)\1_{\mathcal E}(\omega_i)\right)
 \leq C_\vartheta D\varepsilon^{1/3}|A|,
\end{equation}
provided the coefficient in the last count estimate is at most one.
For a single neighborhood the probability of even one exceptional
visitor is bounded by the exponential moment minus one; it is
$O(D\varepsilon^{1/3})$.  With $\varepsilon=CD^{-40}$ this is
smaller than $D^{-6}$.  For a joint bound, choose a fixed large
$\vartheta$ in~\eqref{lp:visitor-mgf} first, and then take $D$ large.
The probability of an exceptional visitor at every site in $A$ is
at most $e^{-(\vartheta-o(1))|A|}$.

To obtain the interior supply, start with the lower endpoint count
$\ell_*D$ supplied by Proposition~\ref{set:low-counts}.
Its failure probability is $O(D^{-2})$, jointly $(C/D^2)^{|A|}$.
Remove labels with a close endpoint at either time, and labels whose
displacement exceeds one.  Proposition~\ref{set:close-tail}
gives a joint small-parameter bound for a fixed positive fraction of
close labels, since $s\ll D^{-1/3}$.  Its one-site probability here
is sharper: by~\eqref{lp:close-pairs} and Markov's inequality, removing
a fixed fraction of $D$ labels has probability
\begin{equation}\label{lp:close-loss}
 CD s^3=CD^{-7/5}.
\end{equation}
Labels initially at a site but with a nonseparated terminal endpoint
are covered by the same estimate in a fixed neighborhood, except
for paths of displacement greater than one.  The latter have free
probability $Ce^{-c/t}$; apply~\eqref{lp:visitor-mgf} to the event
that a fixed positive fraction of the local labels has this property.
Its probability is $e^{-c'D}$ for large $D$.  The stronger cutoff
failures have already been treated.  A fixed positive fraction of
$D$ eligible labels consequently remains.  Conditional on all paths,
independent flags leave at least $cf_0D$ labels with probability
$1-e^{-c'f_0D}$.  The lists based at distinct sites are disjoint, so
the same binomial estimate is joint over any prescribed sites.

We have controlled the supply of eligible labels and the exceptional
visitors.  The remaining certificate tests prevent too many obstacles
from concentrating in small balls.  The upper counts in unit cells
follow from the count exponential moment; the weighted bounds require
a uniform argument down to the shrinking hard-core scale.  Test separated sites, so
far apart that a path obeying the motion cutoff can contribute to
only one of their neighborhoods.  Select one mesh time and one ball
of radius $h_v\in[r,1]$ at each tested site $v$.  Truncate the
oscillation weight at the cutoff.  Multiply each contribution by the
indicator that its whole central-slab path satisfies the increment and
diameter cutoffs; outside the exceptional-visitor event already estimated,
these restricted counts equal the actual local counts.  This is a
restriction of an integrand, not a conditioning of the path law.  Put
$a_D=c'(\log D)^{-2}$.  Since $A_J\leq C\sqrt{\log D}$ on that
cutoff, $a_D(1+A_J)^3=o(1)$.  For the single-path exponential increment
$F_v$ we therefore have
\[
 0\leq F_v\leq C a_D(1+A_J)^3
     \1_{\{\omega(q_v)\in B_v\}},\qquad
 \E_{\mathrm{free},[-1,1]}^y F_v
       \leq C a_D h_v^3e^{-3d_K(y,v)}.
\]
For a following interval the second inequality uses the bounded
Gaussian density at $q_v$ and the bounded moments of the normalized
increment.  For a preceding interval, integrate first over the
increment and translate the ball by that increment; its Gaussian
tail gives the identical bound.  Central-slab times are uniformly
away from the longer slab's endpoints, so the constants are uniform.
For flagged counts, first average the flag of this label; the bound
acquires a factor $f_0$.

Because a path passing the diameter cutoff visits at most one tested
neighborhood, its exponential increment is the sum of the corresponding
$F_v$.  This property survives time reversal: the increment and diameter
gates are reversal invariant, and $F_v\circ\mathrm{rev}$ still requires
a visit to the same spatial neighborhood $v$, merely at the reflected
time and interval.  Thus replacing the single-path increment by
$\sum_v(F_v+F_v\circ\mathrm{rev})$ makes it reversal invariant while
preserving the same separated support and free expectation bound, up to
a factor two.  Expand the product over distinct labels, apply
\eqref{lp:deletion} to these weights, and bound the resulting product
by an exponential.  This gives
\begin{equation}\label{lp:fine-mgf}
 \log\E e^{a_D\sum_v W_v}
 \leq C D a_D\sum_v h_v^3,
\end{equation}
where $W_v$ is the truncated weighted count restricted to paths
satisfying the motion cutoffs.  In the flagged version
the right side has an extra $f_0$.  To justify the use of the spatial
count estimate with different $h_v$, write
$h_v^3=\int_0^1\1_{\{h_v^3\geq u\}}\dd u$.
For $A_u=\{v:h_v^3\geq u\}$, the lattice exponential sum satisfies
$\sum_{v\in A_u}e^{-3d_K(y,v)}\leq C e^{-d_K(y,A_u)}$.
Convexity of the logarithm of the exponential moment, followed by
the set-count bound and
$\int_0^1|A_u|\dd u=\sum_vh_v^3$, proves~\eqref{lp:fine-mgf}.

Exponential Markov inequality at thresholds
$C(Dh_v^3+D^{1/100})$ now bounds any specified collection of ball
violations by
\[
 \exp\{-c a_D D^{1/100}|A|\}.
\]
The same proof applies to flagged counts at their stated thresholds.
There are only polynomially many net balls and mesh times per site;
their union therefore still has an arbitrarily large power saving,
jointly as well as marginally.  A fixed finite coloring separates
arbitrary sites.  Finally assign to every bad site one of the finitely
many failed tests and select a largest class.  The coloring and this
assignment cost only fixed roots and fixed exponential factors in
the joint parameter.  Each underlying parameter can be made as small
as needed before choosing $D$, proving the joint assertion for any
fixed $p$.  Marginally the largest remaining term is
$CD^{-7/5}$ from~\eqref{lp:close-loss}, which is bounded by
$CD^{-1-1/8}$.  Applying the weighted-ball and motion arguments to
all labels instead of only frozen labels proves the weighted part of
the full-path bound.  Either failure in \eqref{lp:ordinary-motion} has
free probability at most $CD^{-40}$ by the first Gaussian estimate and
its reversal.  The visitor estimate consequently bounds their local
union by $O(D\,D^{-40/3})$, which is smaller than $D^{-6}$.  This
proves the asserted full-path failure bound including its ordinary-motion
requirements.
\end{proof}

We have now obtained many labels and quantitative control of the
obstacles in their neighborhoods.  The remaining task is deterministic
in the exposed data: show that a free trial bridge has a large chance
to avoid those obstacles.  The weighted oscillations in
\eqref{lp:frozen-count} were included precisely for this estimate.

\subsection{Collision probabilities and successful insertions}

For distinct labels define
\begin{equation}\label{lp:beta}
 \beta_{ik}=b_D\left(1+
       \frac1{s\vee d_K(Y_i,Y_k)}+
       \frac1{s\vee d_K(Z_i,Z_k)}\right),
 \qquad b_D=\frac{C(\log D)^C}{D}.
\end{equation}
The constants in this polylogarithmic expression are fixed, and may
be enlarged without further mention.

\begin{lemma}[Collision with one allowed path]\label{lp:collision}
Fix the exposed data in a good neighborhood.  Let $k$ be a deterministic
path satisfying its individual cutoffs, and let $i\ne k$ have the
eligible endpoints.  Under $B_i$, the probability that the trial
satisfies its own cutoffs and comes within distance $r$ of path $k$
is at most $\beta_{ik}$.  For two legs pinned at a prescribed midpoint
$y$, with $|y-m_i|\leq8b$ and
$d_K(y,\omega_k(0))>s$, the bound has the additional term
\[
 \frac{b_D}{s\vee d_K(y,\omega_k(0))}.
\]
If $y$ is instead uniform in a unit cell at a fixed bounded distance
from $m_i$, the integrated collision probability, including failure
of this midpoint separation, is at most $C\beta_{ik}$.
\end{lemma}

\begin{proof}
Orient each half of an ordinary bridge from its nearer endpoint;
for a pinned trial do this separately on each of its two legs.  Let
$h\geq s$ be the separation between the trial and opponent at that
endpoint, and let $q$ be elapsed time from it.  On the two increment
cutoffs, a collision in a mesh interval is impossible unless
\[
 q+C\Delta\geq c h^2/(\log D).
\]
The lower bound exceeds $\Delta$ by a power of $D$, since
$s^2/(r^2\log D)\to\infty$.  A possible collision requires positions
within $Cr\sqrt{\log D}$ at a neighboring mesh point.  On the
oriented half-leg, the bridge density there is bounded by $Cq^{-3/2}$.
Thus the sum over these mesh points is at most
\[
 C r^3(\log D)^{3/2}
   \sum_{q\gtrsim h^2/\log D}q^{-3/2}
 \leq \frac{Cr(\log D)^2}{h}.
\]
Summing the finitely many oriented pieces proves the stated endpoint
and midpoint terms, after enlarging $b_D$.  A fixed additive term
covers the harmless bounded-distance conventions.  Confinement
eliminates opponents outside the relevant neighborhood, and local
lifts apply for all sufficiently large $K$.  Finally
$\int_Q(s\vee|y-z|)^{-1}\dd y\leq C$ for every unit cube $Q$ and
every $z$, whereas the volume where $|y-z|\leq s$ is at most $Cs^3$.
Since $s^3=o(b_D)$, integration proves the last assertion.
\end{proof}

The individually small collision bounds sum to $o(1)$ over the
flagged opponents, because their density is $f_0D$.  The frozen
opponents have density of order $D$; for them we shall use the
weighted space-time counts in the certificate.  This distinction gives
the following success bounds, uniformly in all paths left unexposed.

\begin{proposition}[Uniform insertion probabilities]\label{lp:insertion}
Fix good data in all neighborhoods met by the following trials.  In the
first three assertions existing unexposed paths may take any values
satisfying their individual cutoffs, while the frozen paths remain those
of the fixed data.  The fourth assertion instead concerns an obstacle
collection satisfying the full-path bounds.  For all sufficiently large
$D$, these assertions hold also when an arbitrary subset of avoidance
checks is omitted.
\begin{enumerate}
\item An ordinary trial $B_i$, required to satisfy its own cutoffs and
avoid the frozen paths and all indicated other unexposed paths, succeeds
with probability at least $3/4$.
\item A trial formed by choosing a uniform midpoint in a unit cell
contained in $B(m_i,8b)$ and at an absolute constant distance from
its mean, then choosing independent bridge legs, has success probability
at least $3/4$.  It may also be
required to have midpoint distance greater than $s$ from every opponent.
\item Two such trials in two copies, using a common unit cell meeting
the preceding conditions for both labels, with a common uniform midpoint
and independent legs conditional on that midpoint, simultaneously succeed
with probability at least $1/2$.  Each is checked against the previously
present paths in its own copy; the pair is not checked against itself
across the copies.
\item Suppose the entire obstacle collection satisfies the motion and
count bounds defining the full-path event.  A trial pinned at any point
within $8b$ of its mean and more than $s$ from every obstacle midpoint
succeeds with probability at least $3/4$.  This remains valid after adding
or substituting a fixed number of individually allowed obstacles, provided
the indicated endpoint and midpoint separations hold.
\end{enumerate}
\end{proposition}

\begin{proof}
For unexposed opponents, sum Lemma~\ref{lp:collision} over endpoint
shells using~\eqref{lp:flagged-count}.  Both endpoint lists satisfy
\[
 \sum_k\beta_{ik}
 \leq C b_D(f_0D+D^{1/100}/s)
 \leq C(\log D)^C(D^{-x}+D^{-19/100})=o(1).
\]
This also bounds the integrated cost for a uniform midpoint.  The
extra midpoint separation costs at most $Cs^3$ per opponent, and
there are $O(f_0D)$ such opponents in a relevant neighborhood.

For frozen opponents the endpoint-shell bound alone would lose a
logarithm.  We instead use their weighted space-time count.  Consider
an oriented half-leg and put
\[
 q_0=cs^2/(\log D).
\]
Endpoint separation and~\eqref{lp:cutoffs} imply that only intervals
with elapsed time $q\geq q_0$ can contribute.  At their initial mesh
point the trial position has Gaussian density at most
$Cq^{-3/2}e^{-c|z-a_q|^2/q}$, where $a_q$ is its bridge mean.
Given that position, the following increment has the representation
\[
 X_{q+v}=X_q+\frac{v}{\ell-q}(a'-X_q)+R_v,
 \qquad 0\leq v\leq\Delta,
\]
where $\ell\in\{t,2t\}$ is the bridge or leg duration, $a'$ is its
opposite prescribed endpoint, and $\ell-q\geq\ell/2$ on the oriented
half.  The centered bridge $R$ is independent of $X_q$.  Its supremum
divided by $r$, denoted by $U$, has a uniform Gaussian tail.
On confinement the drift
over the mesh interval is at most $C_b\Delta/t\leq r$ for large $D$.
A collision with frozen path $k$ therefore requires
\begin{equation}\label{lp:mesh-collision}
 |\omega_i(q)-\omega_k(q)|
 \leq Cr(1+A_J(\omega_k)+U).
\end{equation}
After deriving this necessary event we discard the cutoff before
integrating the unrestricted bridge, under which the stated independence
holds.  This conditional-increment argument is the one used in
\cite[Proposition~5.2]{H}.

Truncate $U$ at $C\sqrt{\log D}$ with an arbitrarily small polynomial
error, summed over the polynomial number of local paths and mesh
intervals.  The radii in~\eqref{lp:mesh-collision} are then
$o(\sqrt q)$ uniformly for $q\geq q_0$.  Integrate the Gaussian
density over these balls, then integrate $U$.  Its fixed moments are
bounded, so the result is bounded by
\[
 Cr^3q^{-3/2}\sum_k(1+A_J(\omega_k))^3
            e^{-c|\omega_k(q)-a_q|^2/q}.
\]
Gaussian shells of radius $\sqrt q$, and~\eqref{lp:frozen-count},
bound this expression by
$Cr^3(D+D^{1/100}q^{-3/2})$.  Shells of radius greater than one
use a unit-ball covering; their Gaussian decay makes the sum finite.
Summing the mesh intervals on all oriented pieces gives
\begin{equation}\label{lp:frozen-collision-sum}
 C r\left(Dt+\frac{D^{1/100}}{\sqrt{q_0}}\right)+o(1)
 \leq \frac{C\alpha}{x\log D}
      +C\alpha D^{-19/100}\sqrt{\log D}+o(1)=o(1).
\end{equation}
For an ordinary bridge no midpoint is pinned, hence no midpoint
separation is needed.  For a uniform midpoint the excluded volume
around frozen midpoint positions is at most $CDs^3=o(1)$, by the
unweighted count bound.  For each chosen midpoint satisfying the
separation condition, the preceding frozen estimate is uniform.
The individual trial cutoffs fail with probability $O(D^{-40})$.
These estimates prove the first two assertions, and a union bound
in the common-midpoint proposal proves the third.

On the full-path event the entire obstacle collection obeys
\eqref{lp:frozen-count}; the same calculation proves the fourth
assertion for any separated prescribed midpoint.  Each added or
substituted allowed obstacle contributes at most
$Cb_D/s=o(1)$ by Lemma~\ref{lp:collision}.  Deleting checks only
increases each success probability.
\end{proof}

\subsection{Densities paid by a local change}

The midpoint density of $B_i$ in the selected lift is
\[
 g_i(y)=(2\pi t)^{-3/2}
              \exp\{-|y-m_i|^2/(2t)\}.
\]
Indeed Brownian motion has generator $\Delta$, so its duration-$2t$
bridge has midpoint covariance $t$ times the identity.  If a trial
instead has a uniform midpoint in a unit cube $Q$ at an absolute
constant distance from $m_i$, its density relative to $B_i$ is
$\1_Q(y)/g_i(y)$.  Hence that density is at most $CD^{Cx}$.
Conditioning on an insertion event from Proposition~\ref{lp:insertion}
costs only a fixed additional factor.  For a pair sharing its midpoint,
this statement applies to each path marginal, or to an integral that
retains the shared midpoint as an auxiliary variable.  The pair itself
is supported on equality of midpoints and has no density with respect
to the independent product $B_i\otimes B_k$.  When
$|m_i-m_k|\leq C$ and $|y-m_i|\leq Cb$, direct subtraction of the
two Gaussian exponents gives
\begin{equation}\label{lp:midpoint-ratio}
 \frac{g_k(y)}{g_i(y)}
 \leq \exp\!\left\{\frac{C(1+b)}t\right\}
 \leq C D^{C\sqrt{x}}.
\end{equation}
The constants in these exponents are independent of sufficiently
small $x$; only the allowed large-$D$ threshold depends on $x$.

\begin{corollary}[Conditional density domination]\label{lp:trial-density}
Fix good data and a set of route labels $S\subset\mathcal I$ with a bounded
number of labels per site.  In a product of their free bridge
probabilities, condition on individual cutoffs and any selection of
avoidance checks, including checks against fixed outside paths.  The
paths with labels in $\mathcal I\setminus S$ need only satisfy their individual
cutoffs; the frozen paths remain those of the fixed good data.  Every specified $q$
labels have joint marginal density at most $C^q$ relative to the
product of their $B_i$.  The conclusion also holds conditional on
the values of all other labels.  A uniform-midpoint insertion has
the additional density expense $CD^{Cx}$ for each inserted marginal,
or in an augmented integral retaining the common midpoint, and
replacing one free midpoint density by that of a neighboring label
has expense~\eqref{lp:midpoint-ratio}.
\end{corollary}

\begin{proof}
For a single label, its conditional density is an admissibility
indicator divided by its free insertion probability.  Proposition
\ref{lp:insertion} bounds that denominator below by $3/4$, uniformly
in the other allowed paths.  Successively integrate labels, or
equivalently compare the partition integrals before and after each
deletion.  Each deletion costs at most $4/3$, proving $C^q$ for any
specified collection.  The midpoint statements were computed above.
All arguments continue to apply with checks removed because the same
lower insertion bound remains valid.
\end{proof}

The estimates in this section are local in space and uniform in volume.
They supply both the small excluded mass and the bounded conditional
densities needed by the next comparison.  The later use of spatial
routes belongs to the path-intersection approach of
Benjamini, Pemantle, and Peres~\cite{BPP}; the estimates here provide
the exact hard-core path laws on which that route comparison acts.
